\textbf{Phần 1. Trắc nghiệm lựa chọn}
\setcounterpageref{ex}{0}

\begin{ex} %[3H1N07]
Độ dài $1\text{ m}$ bằng bao nhiêu đề-xi-mét?
\choice
{$100\text{ dm}$}
{\True $10\text{ dm}$}
{$1000\text{ dm}$}
{$1\text{ dm}$}
\loigiai{
$1\text{ m} = 10\text{ dm}$.\\
\textbf{Đáp án của câu hỏi là B}
}
\end{ex}

\begin{ex} %[3H1N07]
Độ dài $1\text{ m}$ bằng bao nhiêu xăng-ti-mét?
\choice
{$10\text{ cm}$}
{$1000\text{ cm}$}
{\True $100\text{ cm}$}
{$50\text{ cm}$}
\loigiai{
$1\text{ m} = 100\text{ cm}$.\\
\textbf{Đáp án của câu hỏi là C}
}
\end{ex}

\begin{ex} %[3H1H07]
Đổi $3\text{ m } 4\text{ dm} = \dots\text{ dm}$. Số thích hợp điền vào chỗ chấm là
\choice
{\True $34$}
{$304$}
{$340$}
{$43$}
\loigiai{
$3\text{ m } 4\text{ dm} = 30 + 4 = 34\text{ dm}$.\\
\textbf{Đáp án của câu hỏi là A}
}
\end{ex}

\begin{ex} %[3H1H07]
Hình tam giác có độ dài 3 cạnh là $5\text{ cm}, 6\text{ cm}, 7\text{ cm}$. Chu vi hình tam giác đó là
\choice
{$17\text{ cm}$}
{\True $18\text{ cm}$}
{$19\text{ cm}$}
{$20\text{ cm}$}
\loigiai{
Chu vi: $5 + 6 + 7 = 18\text{ cm}$.\\
\textbf{Đáp án của câu hỏi là B}
}
\end{ex}

\begin{ex} %[3H1H07]
Khối lượng $1\text{ kg}$ bằng bao nhiêu gam? (ở lớp 2 đã làm quen hoặc quy đổi)
\choice
{\True $1000\text{ g}$}
{$100\text{ g}$}
{$10\text{ g}$}
{$500\text{ g}$}
\loigiai{
$1\text{ kg} = 1000\text{ g}$.\\
\textbf{Đáp án của câu hỏi là A}
}
\end{ex}

\begin{ex} %[3H1H07]
Hình tứ giác có độ dài các cạnh lần lượt là $10\text{ cm}, 12\text{ cm}, 15\text{ cm}, 13\text{ cm}$. Chu vi của tứ giác đó là
\choice
{$48\text{ cm}$}
{$52\text{ cm}$}
{$45\text{ cm}$}
{\True $50\text{ cm}$}
\loigiai{
$10 + 12 + 15 + 13 = 50\text{ cm}$.\\
\textbf{Đáp án của câu hỏi là D}
}
\end{ex}

\begin{ex} %[3H1V07]
Độ dài đường gấp khúc $ABCD$ gồm 3 đoạn thẳng: $AB = 12\text{ cm}, BC = 15\text{ cm}, CD = 18\text{ cm}$ là
\choice
{$43\text{ cm}$}
{\True $45\text{ cm}$}
{$47\text{ cm}$}
{$40\text{ cm}$}
\loigiai{
$12 + 15 + 18 = 45\text{ cm}$.\\
\textbf{Đáp án của câu hỏi là B}
}
\end{ex}

\begin{ex} %[3H1V07]
Một can đựng $10\text{ l}$ nước ngọt. Người ta đã rót ra 3 chai, mỗi chai $2\text{ l}$. Lượng nước ngọt còn lại trong can là
\choice
{$3\text{ l}$}
{$5\text{ l}$}
{\True $4\text{ l}$}
{$6\text{ l}$}
\loigiai{
Đã rót: $2 \times 3 = 6\text{ l}$.\\
Còn lại: $10 - 6 = 4\text{ l}$.\\
\textbf{Đáp án của câu hỏi là C}
}
\end{ex}

\begin{ex} %[3H1V07]
Quả dưa hấu cân nặng $3\text{ kg}$, quả bưởi nhẹ hơn quả dưa hấu $1\text{ kg}$. Cả hai quả cân nặng tất cả là
\choice
{\True $5\text{ kg}$}
{$4\text{ kg}$}
{$6\text{ kg}$}
{$7\text{ kg}$}
\loigiai{
Quả bưởi: $3 - 1 = 2\text{ kg}$. Cả hai quả: $3 + 2 = 5\text{ kg}$.\\
\textbf{Đáp án của câu hỏi là A}
}
\end{ex}

\begin{ex} %[3H1V07]
Một hình chữ nhật có chiều dài $15\text{ cm}$, chiều rộng $10\text{ cm}$. Chu vi hình chữ nhật đó là
\choice
{$25\text{ cm}$}
{$40\text{ cm}$}
{$30\text{ cm}$}
{\True $50\text{ cm}$}
\loigiai{
$(15 + 10) \times 2 = 50\text{ cm}$.\\
\textbf{Đáp án của câu hỏi là D}
}
\end{ex}

\begin{ex} %[3H1V07]
Bác thợ may dùng một tấm vải dài $8\text{ m}$ để may 4 bộ quần áo như nhau. Hỏi mỗi bộ quần áo may hết bao nhiêu mét vải?
\choice
{$1\text{ m}$}
{$3\text{ m}$}
{\True $2\text{ m}$}
{$4\text{ m}$}
\loigiai{
$8 : 4 = 2\text{ m}$.\\
\textbf{Đáp án của câu hỏi là C}
}
\end{ex}

\begin{ex} %[3H1C07]
Một đường gấp khúc gồm 4 đoạn thẳng có độ dài bằng nhau. Biết tổng độ dài đường gấp khúc là $36\text{ cm}$. Độ dài mỗi đoạn thẳng là
\choice
{$8\text{ cm}$}
{$7\text{ cm}$}
{$6\text{ cm}$}
{\True $9\text{ cm}$}
\loigiai{
$36 : 4 = 9\text{ cm}$.\\
\textbf{Đáp án của câu hỏi là D}
}
\end{ex}

\textbf{Phần 2. Trắc nghiệm đúng sai}
\setcounterpageref{ex}{0}

\begin{ex} %[3H1H07]
Xét tính đúng sai về các đơn vị đo:
\choiceTF
{\True $1\text{ m} = 100\text{ cm}$.}
{\True $1\text{ dm} = 10\text{ cm}$.}
{$2\text{ m } 5\text{ cm} = 25\text{ cm}$.}
{\True $2\text{ m } 5\text{ cm} = 205\text{ cm}$.}
\loigiai{
\begin{itemize}
    \item $1\text{ m} = 100\text{ cm}$: Đúng.
    \item $1\text{ dm} = 10\text{ cm}$: Đúng.
    \item $2\text{ m } 5\text{ cm} = 205\text{ cm} \neq 25\text{ cm}$. Ý 3 Sai, Ý 4 Đúng.
\end{itemize}
\begin{itemize}
    \item a) \textbf{Đúng}.
    \item b) \textbf{Đúng}.
    \item c) \textbf{Sai}.
    \item d) \textbf{Đúng}.
\end{itemize}
}
\end{ex}

\begin{ex} %[3H1V07]
Một hình tam giác $ABC$ có $AB = 10\text{ cm}, BC = 12\text{ cm}, CA = 15\text{ cm}$. Xét tính đúng sai:
\choiceTF
{\True Cạnh dài nhất là cạnh $CA$ ($15\text{ cm}$).}
{\True Cạnh ngắn nhất là cạnh $AB$ ($10\text{ cm}$).}
{Chu vi hình tam giác $ABC$ là $35\text{ cm}$.}
{\True Chu vi hình tam giác $ABC$ là $37\text{ cm}$.}
\loigiai{
\begin{itemize}
    \item Cạnh dài nhất $CA = 15\text{ cm}$: Đúng.
    \item Cạnh ngắn nhất $AB = 10\text{ cm}$: Đúng.
    \item Chu vi: $10 + 12 + 15 = 37\text{ cm} \neq 35\text{ cm}$. Ý 3 Sai, Ý 4 Đúng.
\end{itemize}
\begin{itemize}
    \item a) \textbf{Đúng}.
    \item b) \textbf{Đúng}.
    \item c) \textbf{Sai}.
    \item d) \textbf{Đúng}.
\end{itemize}
}
\end{ex}

\textbf{Phần 3. Trả lời ngắn}
\setcounterpageref{ex}{0}

\begin{ex} %[3H1H07]
Đổi $4\text{ m } 5\text{ dm}$ ra đề-xi-mét.
\shortans{45}
\loigiai{
\textbf{Đáp số:} $45$\\
$40 + 5 = 45\text{ dm}$.
}
\end{ex}

\begin{ex} %[3H1H07]
Tính chu vi hình tam giác có các cạnh dài $8\text{ cm}, 10\text{ cm}, 12\text{ cm}$.
\shortans{30}
\loigiai{
\textbf{Đáp số:} $30$\\
$8 + 10 + 12 = 30\text{ cm}$.
}
\end{ex}

\begin{ex} %[3H1V07]
Tính chu vi hình vuông có cạnh dài $7\text{ cm}$.
\shortans{28}
\loigiai{
\textbf{Đáp số:} $28$\\
$7 \times 4 = 28\text{ cm}$.
}
\end{ex}

\begin{ex} %[3H1V07]
Đường gấp khúc gồm 3 đoạn thẳng dài lần lượt $14\text{ cm}, 16\text{ cm}, 20\text{ cm}$. Tính độ dài đường gấp khúc theo xăng-ti-mét.
\shortans{50}
\loigiai{
\textbf{Đáp số:} $50$\\
$14 + 16 + 20 = 50\text{ cm}$.
}
\end{ex}

\begin{ex} %[3H1V07]
Có 4 túi gạo, mỗi túi nặng $5\text{ kg}$. Tính tổng khối lượng gạo theo ki-lô-gam.
\shortans{20}
\loigiai{
\textbf{Đáp số:} $20$\\
$5 \times 4 = 20\text{ kg}$.
}
\end{ex}

\begin{ex} %[3H1C07]
Một hình chữ nhật có chu vi $30\text{ cm}$, chiều dài là $9\text{ cm}$. Tính chiều rộng của hình chữ nhật theo xăng-ti-mét.
\shortans{6}
\loigiai{
\textbf{Đáp số:} $6$\\
Nửa chu vi: $30 : 2 = 15\text{ cm}$. Chiều rộng: $15 - 9 = 6\text{ cm}$.
}
\end{ex}
